% Section 17.7, translated from sv/chapters/17/exercises.tex.

\section{Exercises}
\label{c17:sec:uppgifter}

\begin{aside}
\textbf{How to work with the exercises.} Parts~1 and~2 are done during the lesson: first on your
own, a few minutes per exercise, then together. Part~3 is extra exercises for further practice
after the lesson.

\facitnotis{L17}
\end{aside}

% ------------------------------------------------------------------------------------------
\uppgiftsdel{Part 1}{Review exercises}

\uppgift{1.1}{Vectors in the complex plane}
You have the vectors $a = (2, 1)$, $b = (2, -1)$ and $c = (-3, 4)$. In the exercises below, each
vector $(x, y)$ is to be read as a complex number $z = x + jy$.
\begin{parts}
  \item Write the vectors in complex rectangular form.
  \item Draw the vectors in the complex plane (the $x$-axis is the real part, the $y$-axis the
    imaginary part).
  \item Find the lengths of the vectors, that is, the magnitude of each number.
  \item Find the length (magnitude) of $a + b + c$, that is, $\abs{a + b + c}$.
  \item Find the angles of the vectors.
  \item Find a vector $d$ of length $8$ that points in the opposite direction to $c$.
\end{parts}

\uppgift{1.2}{Rectangular form to Euler form}
A voltage $u(t)$ in an AC circuit can be represented by a phasor $U$, which in rectangular form
is written
\[
  U = 3 - j6\,\text{V}.
\]
\begin{parts}
  \item Draw the phasor $U$ in the complex plane (the $x$-axis is the real part, the $y$-axis the
    imaginary part).
  \item Express the phasor $U$ in Euler form, that is, find the magnitude $\abs{U}$ and the phase
    angle $\delta$ such that $U = \abs{U}e^{j\delta}$.
  \item Assume that the frequency of the voltage is $f = 50\,\text{Hz}$. Find the angular
    frequency $\omega$.
  \item Write $u(t)$ as a function of time, $u(t) = \abs{U}e^{j(\omega t + \delta)}$.
\end{parts}

% ------------------------------------------------------------------------------------------
\needspace{16\baselineskip}
\uppgiftsdel{Part 2}{New material}

\uppgift{2.1}{Adding three currents in an AC circuit}
A node in an electric circuit is fed by three sinusoidal currents:
\begin{align*}
  i_1(t) &= 5\sin(\omega t + 30^{\circ})\,\text{mA}, \\
  i_2(t) &= 3\sin(\omega t - 30^{\circ})\,\text{mA}, \\
  i_3(t) &= 4\sin(\omega t + 120^{\circ})\,\text{mA}.
\end{align*}
The total current $i_{\text{tot}}(t)$ in the circuit is
\[
  i_{\text{tot}}(t) = i_1(t) + i_2(t) + i_3(t).
\]
\begin{parts}
  \item Rewrite the currents $i_1(t)$, $i_2(t)$ and $i_3(t)$ as phasors $I_1$, $I_2$ and $I_3$ in
    complex rectangular form.
  \item Calculate the phasor sum $I_{\text{tot}} = I_1 + I_2 + I_3$.
  \item Draw the phasors in the complex plane (the $x$-axis is the real part, the $y$-axis the
    imaginary part).
  \item Convert the result back into a sinusoidal current in the time domain, of the form
    \[
      i_{\text{tot}}(t) = \abs{I_{\text{tot}}}\sin(\omega t + \delta).
    \]
\end{parts}

% ------------------------------------------------------------------------------------------
\uppgiftsdel{Part 3}{Extra exercises}
The exercises below are extra practice, best done on your own after the lesson. Most of them can
be done in your head or with pencil and paper.

\uppgift{3.1}{Phasor from a sinusoidal signal}
Give the phasor in polar form.
\begin{delar}(2)
  \task $u(t) = 12\sin(\omega t + 60^{\circ})$
  \task $u(t) = 5\sin(\omega t - 45^{\circ})$
  \task $i(t) = 2\sin(\omega t)$
  \task $i(t) = 7\sin(\omega t + 180^{\circ})$
\end{delar}

\uppgift{3.2}{Phasors in rectangular form}
Write the phasors in rectangular form.
\begin{delar}(4)
  \task $10\,\angle\,0^{\circ}$
  \task $4\,\angle\,90^{\circ}$
  \task $6\,\angle\,30^{\circ}$
  \task $5\,\angle\,{-53.1^{\circ}}$
\end{delar}

\uppgift{3.3}{Adding two voltages}
Two voltages with the same frequency are $u_1(t) = 3\sin(\omega t)\,\text{V}$ and
$u_2(t) = 4\sin(\omega t + 90^{\circ})\,\text{V}$.
\begin{parts}
  \item Give the phasors in polar form.
  \item Write them in rectangular form.
  \item Calculate the phasor sum and give it in polar form.
  \item Write $u_{\text{tot}}(t)$.
\end{parts}

\uppgift{3.4}{Adding three phasors}
Three voltages are given by $U_1 = 10\,\angle\,0^{\circ}\,\text{V}$,
$U_2 = 5\,\angle\,90^{\circ}\,\text{V}$ and $U_3 = 5\,\angle\,{-90^{\circ}}\,\text{V}$.
\begin{parts}
  \item Write all of them in rectangular form.
  \item Calculate the sum.
  \item Give the sum in polar form.
  \item Explain why $U_2$ and $U_3$ cancel.
\end{parts}

\uppgift{3.5}{Sum and difference of two phasors}
Two voltages are given by $U_1 = 8\,\angle\,30^{\circ}\,\text{V}$ and
$U_2 = 8\,\angle\,{-30^{\circ}}\,\text{V}$.
\begin{parts}
  \item Write both in rectangular form.
  \item Calculate $U_1 + U_2$ and give the answer in polar form.
  \item Calculate $U_1 - U_2$ and give the answer in polar form.
  \item Comment on the results.
\end{parts}

\uppgift{3.6}{Multiplying and dividing phasors}
Calculate.
\begin{delar}(2)
  \task $(5\,\angle\,20^{\circ})(4\,\angle\,40^{\circ})$
  \task $\dfrac{20\,\angle\,100^{\circ}}{5\,\angle\,40^{\circ}}$
  \task $(6\,\angle\,{-30^{\circ}})(2\,\angle\,{-30^{\circ}})$
  \task $\dfrac{9\,\angle\,0^{\circ}}{3\,\angle\,90^{\circ}}$
\end{delar}

\uppgift{3.7}{Impedance of R, L and C}
At $\omega = 2000\,\text{rad/s}$, $R = 100\,\Omega$, $L = 50\,\text{mH}$ and
$C = 5\,\mu\text{F}$.
\begin{parts}
  \item Give $Z_R$ and its phase angle.
  \item Calculate $Z_L = j\omega L$ and give its phase angle.
  \item Calculate $Z_C = -\dfrac{j}{\omega C}$ and give its phase angle.
  \item What do the phase angles mean for the current through each component?
\end{parts}

\uppgift{3.8}{Series RL circuit}
A series circuit has $R = 40\,\Omega$ and $X_L = 30\,\Omega$.
\begin{parts}
  \item Give the impedance $Z$ in rectangular form.
  \item Calculate $\abs{Z}$ and the phase angle.
  \item The circuit is supplied with $U = 100\,\angle\,0^{\circ}\,\text{V}$. Calculate the current.
  \item Does the current lead or lag the voltage?
\end{parts}

\uppgift{3.9}{Series RC circuit}
A series circuit has $R = 60\,\Omega$ and $X_C = 80\,\Omega$.
\begin{parts}
  \item Give the impedance $Z$ in rectangular form.
  \item Calculate $\abs{Z}$ and the phase angle.
  \item The circuit is supplied with $U = 200\,\angle\,0^{\circ}\,\text{V}$. Calculate the current.
  \item Does the current lead or lag the voltage?
\end{parts}

\uppgift{3.10}{Series RLC circuit}
A series circuit has $R = 20\,\Omega$, $X_L = 60\,\Omega$ and $X_C = 45\,\Omega$.
\begin{parts}
  \item Give the impedance $Z$ in rectangular form.
  \item Calculate $\abs{Z}$ and the phase angle.
  \item The circuit is supplied with $U = 50\,\angle\,0^{\circ}\,\text{V}$. Calculate the current.
  \item Is the circuit inductive or capacitive?
\end{parts}

\uppgift{3.11}{Resonance}
A series circuit has $R = 10\,\Omega$, $L = 100\,\text{mH}$ and $C = 10\,\mu\text{F}$. Resonance
occurs at $\omega_0 = \dfrac{1}{\sqrt{LC}}$.
\begin{parts}
  \item Calculate $\omega_0$.
  \item Calculate the resonant frequency $f_0$.
  \item Calculate $X_L$ and $X_C$ at $\omega_0$.
  \item What impedance does the circuit have at resonance?
\end{parts}

\uppgift{3.12}{Voltage division in an AC circuit}
A series circuit with $Z_R = 30\,\Omega$ and $Z_L = j40\,\Omega$ is supplied with
$U = 100\,\angle\,0^{\circ}\,\text{V}$.
\begin{parts}
  \item Calculate $Z_{\text{tot}}$ and $\abs{Z_{\text{tot}}}$.
  \item Calculate the current $I$.
  \item Calculate the voltage $U_R$ across the resistor.
  \item Calculate the voltage $U_L$ across the inductor.
  \item Show that $U_R + U_L = U$ as phasors, even though $\abs{U_R} + \abs{U_L} \neq \abs{U}$.
\end{parts}

\uppgift{3.13}{Finding the phase from an equation}
An AC voltage is given by $u(t) = 10\sin(100\pi t + \delta)\,\text{V}$. At $t = 5\,\text{ms}$,
$u = 5\,\text{V}$.
\begin{parts}
  \item Set up the equation for $\delta$.
  \item Solve the equation as two cases.
  \item Check both solutions.
\end{parts}

\uppgift{3.14}{Impedances in parallel}
The impedances $Z_1 = 10\,\Omega$ and $Z_2 = j10\,\Omega$ are connected in parallel.
\begin{parts}
  \item Calculate $Z = \dfrac{Z_1Z_2}{Z_1 + Z_2}$ in rectangular form.
  \item Calculate $\abs{Z}$ and the phase angle.
  \item Compare with the series connection of the same impedances.
  \item What do the two phase angles have in common?
\end{parts}

\uppgift{3.15}{From the phasor domain to the time domain}
The current $i(t) = 4\sin(500t - 30^{\circ})\,\text{A}$ flows through a circuit. The circuit's
impedance is $Z = 25\,\angle\,60^{\circ}\,\Omega$.
\begin{parts}
  \item Give the phasor of the current.
  \item Calculate the phasor of the voltage, $U = ZI$.
  \item Write $u(t)$.
  \item What is the phase difference of the voltage relative to the current?
\end{parts}

\uppgift{3.16}{Two currents at a node}
A node is fed by $i_1(t) = 6\sin(\omega t + 45^{\circ})\,\text{mA}$ and
$i_2(t) = 8\sin(\omega t - 45^{\circ})\,\text{mA}$.
\begin{parts}
  \item Give the phasors in polar form.
  \item Write them in rectangular form.
  \item Calculate the phasor sum and give it in polar form.
  \item Write $i_{\text{tot}}(t)$.
\end{parts}

\uppgift{3.17}{True or false}
Decide whether each statement is true or false, and justify your answer briefly.
\begin{parts}
  \item Phasors with different angular frequencies may be added.
  \item An inductance has an impedance with the phase angle $+90^{\circ}$.
  \item $\abs{Z} = R + X$
  \item In a purely resistive circuit, current and voltage are in phase.
  \item In a capacitive circuit, the current lags the voltage.
  \item At resonance, the impedance of the circuit is purely resistive.
\end{parts}
